\chapter{Numbers with Units}\label{ch-units}

A number in firmware is almost always a measurement of something: a
delay in milliseconds, a supply in volts, a speed in revolutions per
minute. In most languages the unit lives only in a comment or in the
programmer's head, and mixing two of them up is a classic bug: a delay
of 500 meant as milliseconds and read as microseconds, or millivolts
added to volts. In Zeno a number can carry its unit, and the unit is
checked.

\section{Writing a quantity}

A number followed by a unit is a \emph{quantity}:

\begin{zeno}
let wait = 500 ms
let supply = 3.3 V
let speed = 1200 rpm
print(wait, supply, speed)
\end{zeno}

The units come from a fixed list:

\begin{center}
\begin{tabular}{ll}
\toprule
Measures & Units \\
\midrule
time & \verb|ns us ms s min h| \\
frequency & \verb|Hz kHz MHz| \\
voltage & \verb|uV mV V kV| \\
current & \verb|uA mA A| \\
resistance & \verb|ohm kohm Mohm| \\
power & \verb|mW W kW| \\
energy & \verb|J Wh kWh| \\
temperature & \verb|degC|, and \verb|K| apart from it \\
length & \verb|mm cm m km| \\
mass & \verb|mg g kg| \\
angle & \verb|deg|, and \verb|rad| apart from it \\
rotation & \verb|rpm| \\
speed of a serial line & \verb|baud| \\
\bottomrule
\end{tabular}
\end{center}

A word that is not on the list is refused, with some that are:

\begin{zeno}[fails=yes]
let wait = 5 sec
\end{zeno}

\begin{why}
If any word after a number were a unit, a typing slip would quietly
make a new one, and the bug that units are there to stop would be
back. (Decision 1.15.)
\end{why}

\section{What a quantity can do}

Quantities of the same unit add, subtract and compare as numbers do.
A quantity times or divided by a plain number keeps its unit, and one
quantity divided by another of the same kind is a plain number, how
many times the one goes into the other:

\begin{zeno}
let wait = 500 ms
print(wait + 250 ms, wait * 2, wait / 4, wait > 100 ms)
print(wait / 1 ms, 3.3 V / 1 mV)
\end{zeno}

Dividing by one of a unit, as in \verb|wait / 1 ms|, is how a program
takes the unit off and keeps the bare number, saying which unit the
number is in.

Quantities of different kinds never meet. Volts do not add to seconds,
and a bare number does not add to either:

\begin{zeno}[fails=yes]
let total = 1 s + 3 V
\end{zeno}

\begin{zeno}[fails=yes]
let supply = 3.3 V
print(supply + 2)
\end{zeno}

\section{Units of one kind work together}

Units that measure the same thing combine, and the answer is in the
finer of the two:

\begin{zeno}
print(1 s + 500 ms, 1 s == 1000 ms, 1 h / 1 min, 1 kWh == 3600000 J)
\end{zeno}

\begin{why}
Going from seconds to milliseconds multiplies by a whole number, which
loses nothing. Going the other way divides, and a millisecond is not
exactly a number of seconds in binary. So the direction that is always
exact is the one taken without being asked. (Decision 1.15.)
\end{why}

\subsection{Changing unit}

A quantity is given in another unit of its kind with a dot and the
unit's name. It reads as "delay, in seconds":

\begin{zeno}
let delay = 1500 ms
print(delay.s, (3 s).ms, (3300 mV).V)
\end{zeno}

\section{A unit is a type}

The name of a unit is also a type. A parameter of type \verb|ms|
takes a quantity in milliseconds, and nothing else:

\begin{zeno}
function pause(length: ms) {
    print("pausing for", length)
}
pause(250 ms)
pause((2 s).ms)
\end{zeno}

\begin{zeno}[fails=yes]
function pause(length: ms) {
    print("pausing for", length)
}
pause(2 s)
pause(250)
\end{zeno}

Passing \verb|2 s| where \verb|ms| is wanted is refused, and the
message says to write \verb|.ms|.

\begin{why}
Converting by itself at a call was tried on paper and turned down.
\verb|zeno run| would show \verb|2 s| where a compiled program, which
knows that the parameter is \verb|ms|, would show \verb|2000 ms|: one
program, two outputs. Writing \verb|.ms| means the same in both.
(Decision 1.15.)
\end{why}

\section{Time}

The \verb|time| module works in quantities. \verb|sleep| takes a
time with its unit, and \verb|now()| gives the time in \verb|ms|
since a fixed moment:

\begin{zeno}
from time import sleep, now

let start = now()
sleep(20 ms)
let took = now() - start
print(took >= 20 ms)
\end{zeno}

\verb|sleep(500)| is refused. Whether a bare 500 is seconds or
milliseconds is the question every delay function in every C library
answers differently, and here the program text answers it.

\verb|now()| counts from a moment the program cannot know, such as the
moment a board was switched on, so only the difference of two readings
means anything.

\section{What does not convert}

Degrees Celsius and kelvin differ by an offset, not by a scale, so
adding one to the other means nothing, and Zeno refuses it:

\begin{zeno}[fails=yes]
print(20 degC + 1 K)
\end{zeno}

Degrees and radians of angle have no exact factor between them, so
they too never meet by themselves. A program changes one to the other
when it asks, with \verb|.rad| and \verb|.deg|:

\begin{zeno}
let half_turn = 180.0 deg
print(half_turn.rad, (half_turn.rad).deg)
\end{zeno}

The angle functions of the \verb|math| module take and give angles
with their units (\hyperref[ch-math]{Mathematics}).

\section{Showing a quantity}

A quantity is shown with its unit after a space. A format applies to
the number, and the unit stays:

\begin{zeno}
let supply = 3.3 V
let angle = 30.456 deg
print(f"[{supply:8.2f}] [{angle:.1f}]")
\end{zeno}

The width counts the unit too, so a column of quantities lines up on
the right, as numbers do.

\section{Units made of units}

Two quantities of the right kinds make a third, so Ohm's law, power,
energy and periods are written as they are on paper:

\begin{zeno}
let supply = 3.3 V
let draw = 20 mA
let power = supply * draw
print(power, power.mW)
print(supply / draw, 2 mA * 4.7 kohm)
print(5 W * 2 h, (5 W * 2 h).Wh)
print(1 / 2 ms, 10 ms * 1 kHz)
\end{zeno}

Voltage times current is power, current times resistance is voltage,
and power times time is energy, either way round; each divides back.
A number over a time is a frequency, a number over a frequency is a
time, and a time times a frequency is a plain number, the count of
cycles. The answer is in the unit of its kind with no prefix:
\verb|W|, \verb|A|, \verb|V|, \verb|ohm|, \verb|J|, \verb|s| or
\verb|Hz|, and \verb|.mW| or \verb|.Wh| changes it, as for any
quantity.

Any other pair is refused, with the list of what works:

\begin{zeno}[fails=yes]
let area = 2 m * 3 m
\end{zeno}

\begin{why}
An algebra of dimensions would make \verb|m * m| a unit, and
\verb|V * V| too, and show them as \verb|m^2|. A short list keeps
the rule that every unit has a name, so \verb|V * V|, which is far
more often a slip than a sum, is still caught. (Decision 43.)
\end{why}

There are no units of a program's own yet; that is on the list in
\hyperref[app-limits]{Appendix G}.

\verb|min| and \verb|max| take quantities, as long as all of them are
in one unit, because the answer is one of the values as it was given.
\verb|sum| takes plain numbers only; add quantities with \verb|+|:

\begin{zeno}
let readings = [3.1 V, 2.9 V, 3.4 V]
print(min(readings), max(readings))
var total = 0.0 V
for r in readings {
    total += r
}
print(total)
\end{zeno}

\begin{zeno}[fails=yes]
print(min(1 V, 200 mV))
\end{zeno}

\begin{summary}
A number followed by a unit from Zeno's list is a quantity. Quantities
of one kind combine, in the finer unit, and never mix with another
kind or with bare numbers. \verb|.ms| and its like change unit;
dividing by \verb|1 ms| gives the bare number. A unit's name is a
type. \verb|sleep| and \verb|now| work in quantities. A format applies
to the number, and the unit is always shown.
\end{summary}

\begin{tryit}
\begin{enumerate}
\item Write \verb|function blink_period(rate: Hz) -> ms|. What is the
  bare number you have to work with, and where does the unit go back?
\item Add \verb|250 us| to \verb|1 ms| and print the answer. Then
  print it in microseconds and in seconds.
\item Try \verb|sleep(3 V)| and read the message.
\item Write a voltage divider: given a supply and two resistances in
  \verb|kohm|, give the output voltage as a quantity in \verb|V|.
\end{enumerate}
\end{tryit}
